% !TEX root = ../main.tex

\chapter*{文献与制作说明}
\addcontentsline{toc}{chapter}{文献与制作说明}

\section*{一、底本与框架}

\begin{description}[style=nextline,leftmargin=3.2em]
  \item[汉译底本]
  大正新修大藏经第 99 册《杂阿含经》，求那跋陀罗译（传统称 T99）。
  各经【底本·求那跋陀罗译】即据此出；\textbf{大正经号}（第 1--1362 经）为永久文献锚点。
  \item[卷次与品目]
  本册\textbf{不}按大正五十卷纸面顺序通读，而依 Anesaki／印顺所采用之卷次重排编排「学术卷」。
  每经题下标明学术序与大正经号，便于与馆藏、V1 及既有研究互参。
  \item[底本地位]
  底本用于\textbf{定位经文}与\textbf{对照术语}，并非义理之最终裁判；
  凡与巴利平行明显冲突或文句碍义处，从平行义改写，并于【校勘与说明】标明。
\end{description}

\section*{二、平行经与参考译文}

\begin{description}[style=nextline,leftmargin=3.2em]
  \item[平行关系]
  汉译《杂阿含》与巴利《相应部》等之对应，主要依据
  \textbf{SuttaCentral} 所整理之 curated 平行表（非本书臆测配对）。
  各经【巴利平行与参考译文】列示所据《相应部》等编号（如 22.12）。
  \item[义理裁判]
  译义以巴利平行所显\textbf{早期教理}为主。
  制作过程中参看 SuttaCentral 所收英译以核义；\textbf{书中实际刊出}者见下节「版权与引用」。
  \item[无平行时]
  部分经目在平行表中暂无专经对应（多见于末段或汉本独存经）；
  此类经依汉本与早期阿含教理常识\textbf{保守}处理，【校勘与说明】中会注明依据有限。
\end{description}

\section*{三、本书各层如何产生}

\begin{enumerate}[label=\textbf{\arabic*.}]
  \item \textbf{定位}：依大正经号确定经文；载入汉译底本全文。
  \item \textbf{核义}：查 SuttaCentral 平行及巴利／英译，厘清叙事与法义。
  \item \textbf{主文}：在「信 $>$ 达 $>$ 雅」次序下，以仿鸠摩罗什文体译出【正文·仿罗什风】；
  并写与之段数对应的【今译意】。
  \item \textbf{校勘}：凡据平行改正底本理解处，写入【校勘与说明】；重建经（◇）另述补叙依据。
  \item \textbf{审读}：全书 1362 经均经机器校验（主文与今译意不得雷同、重建经须标记等）及人工抽检闭环。
  \item \textbf{编排}：依 Anesaki／印顺卷次生成阅读序（本册）；插入经入附录。译文与 V1 同源，本阶段不因重排另写经文。
\end{enumerate}

\noindent 本书为\textbf{数字人文与人工审读结合}之研究底本，欢迎依底本、平行与开放数据勘误。

\section*{四、重建经（◇）}

216 经汉本仅为「亦如是」「余如前说」等\textbf{略式}（peyāla），
不足独立阅读。此类经据同品巴利平行或同型经结构补出\textbf{最小可读文}，
经题标 ◇，引用时须注明「重建」。1146 经则有较完整底本经叙，属 full 译出。
详见附录「重建经目录」。

\section*{五、本书刻意不做}

\begin{itemize}
  \item \textbf{不}改写或取代 V1《研究译注本（大正卷序）》；两册并行。
  \item \textbf{不}改动大正经号（\texttt{sa\_t99}）；重排仅影响阅读次序。
  \item \textbf{不}声称唯一正确「原经序」或已恢复出土原典。
  \item \textbf{不}引入后期大乘／禅语汇（如如来藏、佛性等）替代早期用语。
  \item \textbf{不}将重建经伪装为求那跋陀罗译全文。
  \item \textbf{不}声称鸠摩罗什历史上曾译《杂阿含》。
\end{itemize}

\section*{六、版权与引用}

本书为研究译注本：【正文·仿罗什风】与【今译意】为本书原创；
其余层次为\textbf{文献对照}，版权依各来源分别处理。
【巴利平行与参考译文】中巴利、英译均仅刊\textbf{600 字以内摘录}，非全文转载。

\subsection*{6.1 各材料版权状况}

\begin{description}[style=nextline,leftmargin=3.2em]
  \item[汉译底本（T99）]
  大正新修大藏经／CBETA 传统佛典；学界通常作文献引用。
  【底本·求那跋陀罗译】全文刊出，供对照。
  \item[巴利原文]
  来自 SuttaCentral Bilara 所收 \textbf{Mahāsaṅgīti} 数字本。
  巴利经典本身为古文献，无现代著作权；书中仅刊摘录。
  \item[《相应部》英译]
  主要来自 Bhikkhu \textbf{Sujato} 英译（SuttaCentral Bilara）。
  依 \textbf{Creative Commons Zero（CC0）} 公共领域声明：可自由引用、改编与商用；礼貌性署名即可。
  \item[《杂阿含》英译（备用）]
  部分经有 Charles \textbf{Patton} 之 SA 英译（CC0）。
  无 SN 平行而需 SA 英译时，书中可能刊 Patton 或 Anālayo 摘录（见下）。
  \item[Anālayo 英译（少量）]
  约数十经在无 SN 英译时，fallback 为 Bhikkhu \textbf{Anālayo} 英译
  （原刊《法鼓佛学学报》等）。
  \textbf{非 CC0}：Anālayo 授权 SuttaCentral 在其\textbf{网站}刊载，\textbf{并不当然}延及本书。
  本书仅刊短摘录、标注译者及原刊出处；\textbf{未另获 Anālayo 书面授权}；正式商业印行前宜另行确认。
  \item[Bhikkhu Bodhi 等]
  \textbf{本书未嵌入} Bodhi（Wisdom Publications）等出版之英译全文；
  其译本仅可能作为 Sujato 等译者之参考，不在书中逐字刊出。
  \item[平行表与卷次重排]
  SuttaCentral curated 平行关系为参考数据；卷次重排依据 Anesaki（1908）、印顺会编及 Anālayo／Bucknell 等讨论；本书据以编排，不声称原创发现。
\end{description}

\subsection*{6.2 书中摘录与署名}

\begin{itemize}
  \item 有 SN 平行时：优先刊 Sujato 英译摘录，并标「Bhikkhu Sujato（CC0）」。
  \item 无 SN 英译时：可能刊 Patton（CC0）或 Anālayo 摘录；Anālayo 处标注原译者及期刊出处。
  \item 巴利摘录标注「Mahāsaṅgīti 数字本」。
\end{itemize}

\subsection*{6.3 建议引用格式}

\noindent 经号以\textbf{大正经号}为准（与馆藏、V1 一致）；可兼标学术序。示例：

\begin{quote}\small
《杂阿含经·仿罗什风新译》研究译注本（经序重排），大正第 1 经，据《相应部》22.12 校正，2026.

《杂阿含经·仿罗什风新译》研究译注本（经序重排），大正第 275 经，据《增支部》8.9 校正，2026.
\end{quote}

\noindent 引用 ◇ 重建经、Anālayo 英译摘录或 CC0 以外材料时，请一并标明层级与译者。
本书版权说明不构成法律意见；商业出版前建议复核 Anālayo 相关摘录。

\section*{七、开放数据与复核}

研究者可在项目仓库查阅与本书同步之结构化数据（1362 经 JSON）、
学术阅读序表（\texttt{data/metadata/v2/}）、机器审计报告及导出脚本，以便逐经复核、比对改版。
书末「数据与版本」列主要文件名。

\vspace{0.8em}
\noindent\textbf{致谢：}
汉译底本承佛典整理传统；巴利语料及 CC0 英译（Bhikkhu Sujato、Charles Patton 等）来自 SuttaCentral Bilara；
Anālayo 英译原刊《法鼓佛学学报》等，由 SuttaCentral 经译者授权在其平台刊载，本书仅刊短摘录并标注译者。
卷次重排假说承 Anesaki、印顺法师及当代相应部／杂阿含研究者。
本研究译本、体例与排版：Agama–Kumarajiva Project，2026.

\clearpage
